\documentclass[aspectratio=169,spanish,xcolor=table]{beamer}
\usepackage{iftex}
\ifPDFTeX
  \usepackage[utf8]{inputenc}
  \usepackage[T1]{fontenc}
\fi
\usepackage[spanish,es-tabla,es-noshorthands]{babel}
\usepackage{helvet}
\renewcommand{\familydefault}{\sfdefault}
\usepackage{tikz,pgfplots,array,booktabs,ragged2e}
\pgfplotsset{compat=1.18}
\usetikzlibrary{arrows.meta,positioning,calc,shapes.geometric,fit}

\definecolor{iiiepeNavy}{HTML}{163B5C}
\definecolor{iiiepeBlue}{HTML}{1F6F9F}
\definecolor{iiiepeAqua}{HTML}{20B7A8}
\definecolor{iiiepeOrange}{HTML}{F28C28}
\definecolor{iiiepePaper}{HTML}{F6F8F7}
\definecolor{iiiepeInk}{HTML}{1E2A32}
\definecolor{iiiepePale}{HTML}{EEF5F8}
\definecolor{iiiepeMint}{HTML}{E7F6F4}
\definecolor{iiiepeWarm}{HTML}{FFF3E6}

\usetheme{default}
\usecolortheme{default}
\setbeamertemplate{navigation symbols}{}
\setbeamercolor{normal text}{fg=iiiepeInk,bg=white}
\setbeamercolor{structure}{fg=iiiepeNavy}
\setbeamercolor{frametitle}{fg=white,bg=iiiepeNavy}
\setbeamercolor{block title}{fg=white,bg=iiiepeBlue}
\setbeamercolor{block body}{fg=iiiepeInk,bg=iiiepePale}
\setbeamercolor{item}{fg=iiiepeOrange}

\setbeamertemplate{frametitle}{%
  \nointerlineskip%
  \begin{beamercolorbox}[wd=\paperwidth,ht=0.95cm,dp=0.22cm,leftskip=0.45cm,rightskip=0.45cm]{frametitle}%
    \begin{minipage}[c]{13.2cm}\usebeamerfont{frametitle}\large\bfseries\insertframetitle\end{minipage}%
    \hfill\begin{minipage}[c]{1.6cm}\raggedleft\includegraphics[height=0.56cm]{img/iiiepe/logo-iiiepe-blanco.png}\end{minipage}%
  \end{beamercolorbox}%
}

\setbeamertemplate{footline}{%
  \leavevmode%
  \hbox{%
    \begin{beamercolorbox}[wd=.78\paperwidth,ht=0.5cm,dp=0.18cm,leftskip=0.45cm]{frametitle}%
      \scriptsize\color{white}M. de la Cruz Muñoz · MAEM · Matrícula 00874 · SD26 · Actividad 5 · Aprendizaje Profundo%
    \end{beamercolorbox}%
    \begin{beamercolorbox}[wd=.22\paperwidth,ht=0.5cm,dp=0.18cm,rightskip=0.45cm]{frametitle}%
      \hfill\scriptsize\color{white}\insertframenumber{} / \inserttotalframenumber%
    \end{beamercolorbox}%
  }%
  \vskip0pt%
}

\newcommand{\Pill}[2]{%
  \tikz[baseline=(x.base)]{\node[rounded corners=3pt,fill=#1!14,draw=#1!75!black,thick,inner xsep=6pt,inner ysep=4pt,align=center] (x) {\scriptsize\bfseries #2};}%
}

\title{Aprendizaje Profundo y Cultura Escolar}
\subtitle{Actividad 5 · Bloque III · Módulo 51707}
\author{Mtro. Martín Jonathan de la Cruz Muñoz · Matrícula 00874}
\institute{IIIEPE · Maestría en Aprendizaje y Enseñanza de las Matemáticas · Grupo A}
\date{19 de noviembre de 2026}

\begin{document}

\begin{frame}[plain]
  \begin{tikzpicture}[remember picture,overlay]
    \fill[iiiepeNavy] (current page.north west) rectangle ([xshift=10.3cm]current page.south west);
    \fill[iiiepeBlue] ([xshift=10.3cm]current page.north west) rectangle (current page.south east);
    \fill[iiiepeOrange] ([xshift=0.8cm,yshift=-6.45cm]current page.north west) rectangle ([xshift=9.3cm,yshift=-6.53cm]current page.north west);
    \node[anchor=north west,text width=8.7cm,align=left] at ([xshift=0.8cm,yshift=-1.15cm]current page.north west) {%
      {\color{iiiepeAqua!45!white}\small\bfseries MAESTRÍA EN APRENDIZAJE Y ENSEÑANZA DE LAS MATEMÁTICAS · SD26\par}%
      \vspace{0.35cm}
      {\color{white}\Large\bfseries \inserttitle\par}%
      \vspace{0.25cm}
      {\color{white!90}\normalsize \insertsubtitle\par}%
    };
    \node[anchor=south west,text width=8.7cm,align=left] at ([xshift=0.8cm,yshift=0.75cm]current page.south west) {%
      {\color{white}\normalsize\bfseries Mtro. Martín Jonathan de la Cruz Muñoz\par}%
      {\color{white!85}\scriptsize Matrícula 00874 · martin.delacruz@campus.iiiepe.edu.mx\par}%
      {\color{iiiepeAqua!35!white}\scriptsize Monterrey, Nuevo León · 19 de noviembre de 2026\par}%
    };
    \node[anchor=center,text width=5.0cm,align=center] at ([xshift=13.15cm,yshift=-4.5cm]current page.north west) {%
      \includegraphics[width=3.6cm]{img/iiiepe/logo-iiiepe-blanco.png}\par\vspace{0.45cm}%
      {\color{white}\large\bfseries IIIEPE\par}%
      \vspace{0.1cm}%
      {\color{white!85}\scriptsize Fundamentos para la enseñanza\\y el aprendizaje I\par}%
    };
  \end{tikzpicture}
\end{frame}
\begin{frame}{Aprendizaje Profundo y Cultura Escolar (Bloque III)}
  \begin{block}{Tesis de Quiroga Aguilar y Lara Órdenes (2022)}
    Superar la enseñanza transmisiva en matemáticas exige transformar no solo las actividades de clase, sino las relaciones entre docente y estudiantado y las creencias que sostienen la cultura escolar.
  \end{block}
  \vspace{0.3cm}
  \centering
  \Pill{iiiepeNavy}{Superficial vs Profundo} \hspace{0.2cm}
  \Pill{iiiepeBlue}{3 Dimensiones} \hspace{0.2cm}
  \Pill{iiiepeAqua!75!black}{Algoritmo vs Sentido} \hspace{0.2cm}
  \Pill{iiiepeOrange!85!black}{Nuevas Pedagogías}
\end{frame}

\begin{frame}{1. Aprendizaje Superficial vs Aprendizaje Profundo}
  \scriptsize\renewcommand{\arraystretch}{1.4}
  \begin{tabular}{p{2.4cm}p{5.3cm}p{5.3cm}}
    \toprule
    \rowcolor{iiiepePale} \textbf{Criterio} & \textbf{Aprendizaje superficial (Marton y Säljö, 1976)} & \textbf{Aprendizaje profundo (Sawyer, 2008; Quiroga y Lara, 2022)} \\
    \midrule
    Propósito & Memorizar definiciones o pasos para cumplir el requisito. & Comprender conceptos, explicar causas y transferir. \\
    Organización & Ideas sueltas sin conexión entre representaciones. & Red articulada entre tabla, gráfica, símbolo y fenómeno. \\
    Papel ante el reto & Dependencia de ejemplos idénticos mostrados en clase. & Exploración crítica y ajuste autónomo de estrategias. \\
    \bottomrule
  \end{tabular}
\end{frame}

\begin{frame}{2. Enseñanza Tradicional de las Matemáticas}
  \begin{block}{Rasgos que inducen aprendizaje superficial}
    \begin{itemize}\small
      \item Exponer reglas y despejes antes de problematizar su sentido físico o lógico.
      \item Saturar con ejercicios mecánicos idénticos que premian velocidad sin comprensión.
      \item Aislar álgebra, geometría y datos de situaciones reales del entorno.
      \item Evaluar exclusivamente el resultado numérico final sin exigir argumentación.
    \end{itemize}
  \end{block}
\end{frame}

\begin{frame}{3. ¿Por Qué el Algoritmo No Basta?}
  \begin{columns}[T]
    \column{0.48\textwidth}
      \begin{block}{Fluidez procedimental (insuficiente)}
        $$0.5x + 2 = 0.2x + 3.8 \implies x = 6$$
        \scriptsize El alumno ejecuta correctamente reglas de despeje lineal sin analizar variables, unidades ni dominio físico.
      \end{block}
    \column{0.48\textwidth}
      \begin{block}{Comprensión profunda (necesaria)}
        \scriptsize Interpretar $y_1=0.5x+2$ frente a $y_2=0.2x+3.8$:
        \begin{itemize}\tiny
          \item Reconocer ordenadas al origen ($2$ y $3.8$) y distinguir función afín de proporcional $y=kx$.
          \item Justificar en qué intervalo conviene un sensor o costo frente al otro antes y después de $x=6$.
        \end{itemize}
      \end{block}
  \end{columns}
\end{frame}

\begin{frame}{4. Triada de Dimensiones del Aprendizaje Profundo}
  \centering
  \begin{tikzpicture}[font=\scriptsize, >=Latex, scale=0.86, transform shape]
    \node[draw=iiiepeNavy,fill=iiiepePale,thick,circle,minimum size=2.0cm,text width=1.8cm,align=center,font=\scriptsize\bfseries,text=iiiepeNavy] (c) at (0,0) {Aprendizaje\\Matemático};
    \node[draw=iiiepeNavy,fill=white,thick,rounded corners=4pt,text width=5.0cm,align=center] (a) at (0,2.35) {\textbf{\color{iiiepeNavy}Cognitiva:} Pensamiento crítico, comprensión relacional y transferencia.};
    \node[draw=iiiepeBlue,fill=white,thick,rounded corners=4pt,text width=4.8cm,align=center] (b) at (-4.3,-1.7) {\textbf{\color{iiiepeBlue}Interpersonal:} Colaboración, escucha y argumentación matemática entre pares.};
    \node[draw=iiiepeAqua!80!black,fill=iiiepeMint,thick,rounded corners=4pt,text width=4.8cm,align=center] (d) at (4.3,-1.7) {\textbf{\color{iiiepeAqua!75!black}Intrapersonal:} Autorregulación, autoeficacia y mentalidad de crecimiento.};
    \draw[<->,very thick,iiiepeNavy] (a) -- (b);
    \draw[<->,very thick,iiiepeAqua!80!black] (a) -- (d);
    \draw[<->,very thick,iiiepeOrange!85!black] (b) -- (d);
  \end{tikzpicture}
\end{frame}

\begin{frame}{5. Transformación de Roles y Tres Planos de Cambio}
  \begin{columns}[T]
    \column{0.50\textwidth}
      \begin{block}{Transición relacional (Fullan y Langworthy, 2014)}
        \begin{itemize}\scriptsize
          \item \textbf{Control docente} $\rightarrow$ \textbf{Asociación} $\rightarrow$ \textbf{Autonomía}.
          \item El profesorado pasa de transmisor único a activador y diseñador de experiencias auténticas.
        \end{itemize}
      \end{block}
    \column{0.46\textwidth}
      \begin{block}{Tres planos en cultura escolar}
        \begin{itemize}\scriptsize
          \item \textbf{Recursos:} Tareas abiertas auténticas.
          \item \textbf{Didáctica:} Andamiaje y colaboración.
          \item \textbf{Creencias:} Qué entendemos por aprender y evaluar matemáticas.
        \end{itemize}
      \end{block}
  \end{columns}
\end{frame}

\begin{frame}{Evaluación Formativa y Experto Adaptativo}
  \begin{itemize}\normalsize
    \item Evaluar en profundidad requiere evidencias de proceso, devolución descriptiva y criterios claros antes de la tarea \textbf{(Collin y Quigley, 2021; Boström y Palm, 2023)}.
    \item El docente actúa como \textbf{experto adaptativo}: observa qué obstáculos persisten, cuestiona sus propias rutinas e indaga entre colegas.
    \item La colaboración escolar cobra fuerza cuando es específica, deliberada y orientada al aprendizaje real del alumnado.
  \end{itemize}
\end{frame}

\begin{frame}{Transferencia en Bachillerato Técnico e Ingeniería}
  \begin{itemize}\normalsize
    \item Sustituir listas mecánicas por casos situados de calibración, costos y tolerancias en taller.
    \item Exigir defensa oral o escrita de unidades, tasas de cambio y dominio físico de validez.
    \item Integrar reflexión metacognitiva individual tras el trabajo colaborativo en equipos.
  \end{itemize}
\end{frame}

\begin{frame}{Conclusiones de la Actividad 5}
  \begin{itemize}\normalsize
    \item El dominio algorítmico aislado no equivale a comprensión profunda ni transferencia matemática.
    \item Solo al integrar lo cognitivo, lo interpersonal y lo intrapersonal cambia la experiencia del aula.
    \item La cultura escolar se transforma cuando se revisan colectivamente las creencias sobre la enseñanza, preparando el diseño de innovación de la Actividad 6.
  \end{itemize}
\end{frame}

\begin{frame}{Referencias Principales (APA 7)}
  \scriptsize\setlength{\parskip}{5pt}
  \begin{itemize}
    \item Boström, E., y Palm, T. (2023). Formative assessment practice in mathematics. \textit{Frontiers in Education}, 8, 1101192.
    \item Collin, J., y Quigley, A. (2021). \textit{Teacher feedback to improve pupil learning}. EEF.
    \item Fullan, M., y Langworthy, M. (2014). \textit{A rich seam: How new pedagogies find deep learning}. Pearson.
    \item Marton, F., y Säljö, R. (1976). On qualitative differences in learning. \textit{Br. J. Educ. Psychol.}, 46(1), 4--11.
    \item Quiroga Aguilar, L., y Lara Órdenes, E. (2022). El aprendizaje profundo como herramienta para cambio en la visión de aprendizaje de una cultura escolar. \textit{REA}, 12(1).
    \item Sawyer, R. K. (2008). Optimizing learning: Implications of learning sciences research. OECD.
    \item Schunk, D. H. (2012). \textit{Teorías del aprendizaje} (6.ª ed.). Pearson.
  \end{itemize}
\end{frame}

\end{document}
